\section{Discussion}

The multi-seed benchmark shows a clear accuracy--deployment trade-off: Random Forest and LSTM reached lower average macro errors, while the SNN remained competitive and clearly more stable than MLP. \rev{Wilcoxon and Friedman outcomes and the overlapping confidence intervals indicate that Random Forest, LSTM, and the SNN are statistically indistinguishable on this dataset, so the SNN should be read as a robust edge-oriented alternative rather than a universal best-accuracy model.}
From the systems perspective, platform behavior is heterogeneous: Device B reached the highest throughput, Device C the lowest incremental energy per inference, and ESP32 full offline feasibility at much lower processing rates. Because baseline power varies across platforms, energy per inference combined with throughput is more actionable than absolute watts. Augmentation ablation showed a modest but consistent macro-MAE improvement (2.53\%), indicating that controlled synthetic expansion stabilizes training in this low-sample regime.

\subsection{Practical Deployment Guidance}
The results yield a simple field guideline: smartphone or SBC-class execution suits dense sampling, frequent map updates, or near-real-time dashboards; ESP32 remains viable for sparse periodic checks despite higher mJ/inf. Model selection should not rely on minimum MAE alone: prefer a multi-objective envelope combining acceptable error, seed stability, throughput, and incremental energy per inference, extensible with battery and telemetry constraints.

\rev{Although Random Forest achieved the lowest mean MAE, its gap to the SNN is not statistically significant (Wilcoxon tests and confidence intervals fully overlap), while the SNN offers properties tree ensembles lack: event-driven computation whose cost scales with spike activity, a single compact model exportable to ONNX/TFLite on all three platform classes, forward compatibility with neuromorphic accelerators, and a tunable accuracy--energy knob through the temporal parameter $T$ --- i.e., the statistically equivalent model with the best deployment envelope.}

\rev{Cloud inference offers larger models and simpler updates, but in remote Amazonian deployments it adds radio energy that typically dominates the per-sample budget, latency tied to intermittent links, recurring costs, and exposure of raw data to third parties. Fully offline inference inverts this profile: energy is bounded by local computation, monitoring continues during outages, and only aggregated results need eventual synchronization; a hybrid local-screening/cloud-consolidation pattern is a natural evolution.}

\subsection{Error Interpretation and Validity}
Predicted-vs-true plots show a target-dependent error profile: pH and OD follow the identity trend with lower dispersion, while turbidity spreads wider, consistent with freshwater heterogeneity. This supports a layered decision strategy: treat pH and OD as short-term screening signals and turbidity mainly as a flag for additional sampling. \rev{For dissolved oxygen specifically, no evaluated model achieves positive mean $R^2$, pointing to a feature-set rather than a modeling limitation: the current inputs carry little OD-related information ($|r| \leq 0.24$), so improving OD prediction requires enriching the input modality (e.g., conductivity, oxidation-reduction potential, chlorophyll, or diel temperature) rather than further model tuning.}

Threats to validity remain centered on sample size and regional representativeness, requiring multi-season, multi-site campaigns with harmonized protocols and external validation. \rev{Transferability should be framed at the pipeline level rather than the trained weights: the learned spatial gradients are specific to Lago do Serpa, so what transfers is the methodology (feature design, augmentation policy, statistical protocol, and deployment stack), retrainable from a comparably small campaign at a new site; cross-site fine-tuning and domain adaptation are planned to quantify this transfer cost.} Reproducibility remains a practical strength: the protocol and measurements are explicit and portable, and sub-millisecond Raspberry Pi inference plus autonomous ESP32 execution confirm practicality across heterogeneous edge profiles.
